% Portability section (\input by paper.tex). Every capability claim is from
% the official documentation page cited next to it, accessed 2026-11-21.

Nothing in \cz{}'s gate, cleaner, audit log or hardware root is specific
to Claude Code: they take a tool name and a JSON input and return a
decision, or take a tool result and return a cleaned one. What is
specific is the \emph{encapsulation} of \cref{sec:inject,sec:session},
and it rests on four capabilities of the host CLI. \textbf{R1}, a
blocking pre-tool hook: an external program (or, at worst, an in-process
plugin) that receives every tool call, in the main loop \emph{and in
subagents}, before it runs, can deny it, and fails closed. \textbf{R2}, a
post-tool output rewrite, which the cleaner needs. \textbf{R3}, a settings
injection channel ranked above what a repository can supply, plus a way
to see and refuse hooks and command-valued settings the project adds.
\textbf{R4}, an on-disk transcript of every tool call with a stable call
identifier, for the binding. R1 is necessary; R2 serves only the cleaner;
R3 and R4 make the monitor \emph{reliable} against a hostile repository
and hostile processes, and without them the corresponding attacks are
detected later or not at all. \Cref{tab:portability} gives what each
CLI's documentation says (``n.d.'': the cited page does not say;
nothing is inferred from source code or third-party material).

\begin{table}[!htbp]
\centering\footnotesize
\caption{Agent CLIs against requirements R1--R4, from their official
documentation (accessed 2026-11-21). Fail: what a failed or timed-out
pre-tool hook does by default.}
\label{tab:portability}
\setlength{\tabcolsep}{3pt}
\begin{tabularx}{\linewidth}{@{}>{\raggedright\arraybackslash}p{0.2\linewidth}>{\raggedright\arraybackslash}X>{\raggedright\arraybackslash}p{0.13\linewidth}>{\raggedright\arraybackslash}p{0.17\linewidth}>{\raggedright\arraybackslash}p{0.12\linewidth}@{}}
\toprule
CLI & R1 pre-tool hook (fail; subagents) & R2 rewrite & R3 injection & R4 transcript \\
\midrule
Claude Code \cite{cc-hooks,cc-settings,cc-cli} & Yes (open; closed with \texttt{onFailure: "block"}; subagents yes) & Yes & Yes (\texttt{-{}-settings}) & Yes \\
Codex CLI \cite{codex-hooks,codex-cli} & Yes, not hosted tools (open; n.d.) & Partial & Partial & Yes (unstable) \\
Gemini CLI \cite{gemini-hooks,gemini-hooks-ref,gemini-config} & Yes (open; n.d.) & Yes & Partial & Yes \\
Cursor \cite{cursor-hooks,cursor-cli} & Yes, in the agent (open unless \texttt{failClosed}; n.d.) & MCP only & Partial & Yes \\
Copilot CLI \cite{copilot-hooks} & Yes (errors deny, timeouts open; implied) & Yes & Partial & Partial \\
Kiro CLI \cite{kiro-hooks,kiro-hook-types,kiro-hook-actions} & Yes (non-zero exit blocks; n.d.) & n.d. & No & n.d. \\
OpenCode \cite{opencode-plugins-v1,opencode-plugins-v2,opencode-config-v1} & In-process plugin (n.d.; a subagent bypass reported \cite{opencode-5894}) & Yes & Yes & n.d. \\
Amp \cite{amp-plugin-api,amp-settings,amp-manual} & In-process plugin (n.d.; no subagent tools) & Yes & Partial & n.d. \\
Cline \cite{cline-plugins} & In-process SDK plugin (errors ignored; n.d.) & Yes & Partial & n.d. \\
Aider \cite{aider-options} & No & No & config file & Partial \\
\bottomrule
\end{tabularx}
\end{table}

Hooks have become common: every CLI surveyed except Aider documents a
pre-tool hook that can deny a call. The gaps are the ones that matter in
Claude Code: most hooks fail \emph{open} by default, subagent coverage is
stated explicitly only for Claude Code, and only Claude Code documents a
command-line settings file ranked above project configuration, so in
most others a launcher must find and refuse the hooks a repository adds.
For a CLI that meets R1--R4 with external-program hooks (Codex CLI and
Gemini CLI come closest), a port is an adapter: translate the CLI's hook
input into \sys{}'s call shape and the decision back, generate its hook
configuration in a private session directory, and parse its transcript
for the binding; the session file, \texttt{env -i} start, control
socket, peer check, handshake, gate, cleaner, audit log and hardware
root carry over unchanged. Where a hook fails open, the launcher must
make failure unlikely and rely on the transcript watch. For in-process
plugins (OpenCode, Amp, Cline), a thin shim forwards each call to the
external hook, and a project plugin loaded into the same process must be
refused at launch. Only the Claude Code port has been built and
measured; the other ports are design, not results.
